\begin{figure}[H]
\centering
\begin{tikzpicture}[font=\small,>=Latex,xscale=1.35,yscale=1.2]
  \fill[blue!6] (-2,-2) rectangle (2,0);
  \fill[orange!22]
    (-1.12,0)--(-1.12,-.35)
    .. controls (-1.12,-.83) and (-.57,-1.28) .. (0,-1.28)
    .. controls (.57,-1.28) and (1.12,-.83) .. (1.12,-.35)
    --(1.12,0)--cycle;
  \draw[thick]
    (-1.12,0)--(-1.12,-.35)
    .. controls (-1.12,-.83) and (-.57,-1.28) .. (0,-1.28)
    .. controls (.57,-1.28) and (1.12,-.83) .. (1.12,-.35)
    --(1.12,0);
  \draw[very thick,blue!65!black] (-2,0)--(-2,-2)--(2,-2)--(2,0);
  \draw[thick,dashed] (-2,0)--(2,0);
  \node[above=7pt] at (0,0) {artificial edge \(s=b\)};
  \node at (-.2,-.48) {\(K<0\)};
  \node at (0,-1.65) {\(K>0\)};
  \node[anchor=east] at (-1.15,-.85) {\(K=0\)};
  \draw[<->] (.66,-.04)--(.66,-.82);
  \node[anchor=west] at (.72,-.58) {\(d\)};
  \draw[->,gray] (2.45,-1.8)--(2.45,-.65) node[above,black] {\(s\)};
  \draw[->,gray] (2.2,-1.55)--(3.05,-1.55) node[right,black] {\(t\)};
  \node[text=blue!65!black] at (0,-2.5) {three edges in positive curvature};
\end{tikzpicture}
\caption{A lower-cap integration region in flow coordinates, schematically
and not to scale. The lateral and bottom edges are chosen in positive
curvature. The artificial edge \(s=b\) may cross the negative region;
no high-order data are prescribed there. The distance \(d=b-s\)
enters the vanishing weight. The actual curvature zero set depends on
the solution through the flow coordinates; only its position relative
to the three elliptic edges and the artificial edge is illustrated.}
\label{fig:cap}
\end{figure}
